
\def \Subject {بهینه‌سازی}
\def \Session { دو}
% \setcounter{chapter}{\Session-1}
\renewcommand{\chaptername}{سوال}
\chapter{\Subject}

\section{پیاده‌سازی تابع هدف و گرادیان}
مطابق خواسته سوال تابع هدف و گرادیان پیاده‌سازی شد. نمودار این تابع در شکل 
\ref{fig:rosenblot}
آمده است. 

\begin{figure}[H]
	\centering
	\includegraphics[width=0.8\textwidth]{images/q2_rosenbrock.png}
	\caption{نمودار تابع \lr{rosenbrock} بین یک تا منفی یک}
	\label{fig:rosenbrock}
\end{figure}

\section{پیاده‌سازی بهینه‌سازها}
مطابق خواست سوال برای توابع بهینه‌سازی هر سه تابع بهینه‌ساز پیاده شدند و مسیرهای آن‌ها در شکل 
\ref{fig:rosenbrock_path}
آمده است.  

\begin{figure}[H]
	\centering
	\includegraphics[width=0.8\textwidth]{images/q2_paths.png}
	\caption{نمودار تابع \lr{rosenbrock} بین یک تا منفی یک}
	\label{fig:rosenbrock_path}
\end{figure}

همچنین نمودارهای فاصله و خطا برای مقادیر خواسته شده حساب شد. جهت اینکه مقادیر قابل مشاهده باشند از آن‌ها لگاریتم گرفته شده است. 

\begin{figure}[H]
	\centering
	\includegraphics[width=0.8\textwidth]{images/q2_distance0.01.png}
	\caption{نمودار فاصله برای ضریب یادگیری یک صدم.}
\end{figure}
\begin{figure}[H]
	\centering
	\includegraphics[width=0.8\textwidth]{images/q2_distance0.1.png}
	\caption{نمودار فاصله برای ضریب یادگیری یک دهم.}
\end{figure}
\begin{figure}[H]
	\centering
	\includegraphics[width=0.8\textwidth]{images/q2_distance1.0.png}
	\caption{نمودار فاصله برای ضریب یادگیری یک .}
		\label{fig:dist_1}
\end{figure}
\begin{figure}[H]
	\centering
	\includegraphics[width=0.8\textwidth]{images/q2_loss0.01.png}
	\caption{نمودار مقدار تابع برای ضریب یادگیری یک صدم.}
\end{figure}
\begin{figure}[H]
	\centering
	\includegraphics[width=0.8\textwidth]{images/q2_loss0.1.png}
	\caption{نمودار مقدار تابع برای ضریب یادگیری یک دهم.}
\end{figure}
\begin{figure}[H]
	\centering
	\includegraphics[width=0.8\textwidth]{images/q2_loss1.0.png}
	\caption{نمودار مقدار تابع برای ضریب یادگیری یک .}
	\label{fig:loss_1}
\end{figure}

همانطور که در شکل 
\ref{fig:loss_1}
و 
\ref{fig:dist_1}
می‌توان دید، هنگامی که نرخ یادگیری زیاد می‌شود، مدل 
\lr{converge}
نمی‌کند. 

همچنین در همه مدل‌ها 
\lr{SGD}
عملکرد خیلی بدی دارد زیرا تابع محدب نیست.